La superfície de la Terra és principalment aigua que conté ions en dissolució i que li fan adquirir una càrrega neta negativa. Es pot considerar que la Terra té un camp elèctric en punts propers a la seva superfície amb un mòdul constant de 150~N/C.

\begin{apartat}{1,25}
Dibuixeu l'esfera terrestre i representeu-hi el camp elèctric al voltant de la superfície. Calculeu el valor de la càrrega total que produeix aquest camp elèctric. Per fer-ho, considereu que el camp elèctric creat per una superfície esfèrica carregada uniformement és igual al generat per tota la càrrega situada al centre de l'esfera.
\end{apartat}

\begin{apartat}{1,25}
Calculeu el mòdul de la força elèctrica que produirà el camp elèctric sobre un electró lliure situat a la vora de la superfície de la Terra. Calculeu la massa que ha de tenir una gota esfèrica d'aigua amb una càrrega extra d'un sol electró perquè el seu pes es compensi amb la força elèctrica. Feu un esquema en què es mostrin les forces que actuen sobre la gota. Calculeu el diàmetre d'aquesta gota d'aigua.
\end{apartat}

\dades{%
  Radi de la Terra, $R_\mathrm{T} = 6,37\times 10^{6}\ \mathrm{m}$.\\
  Densitat de l'aigua, $\rho = 10^{3}\ \mathrm{kg/m^3}$.\\
  Massa de l'electró, $M_\mathrm{e} = 9,11\times 10^{-31}\ \mathrm{kg}$.\\
  Superfície esfèrica: $4\pi r^2$.\\
  Volum d'una esfera: $4/3\pi r^3$.\\
  $|e| = 1,602\times 10^{-19}\ \mathrm{C}$.\\
  $k = \dfrac{1}{4\pi\varepsilon_0} = 8,99\times 10^{9}\ \mathrm{N\,m^2\,C^{-2}}$.\\
  $g = 9,81\ \mathrm{m/s^2}$.}
